\section{The rational primitive and its localization}\label{h3:sec:rational}

The integral basis gives a distinguished line
$\Z_p\delta\subset\pi_{-3}L_{K(1)}T$.  We now construct a rational
class in $\pi_{-3}L_0T$ whose image is nonzero in its rational span
$\Q_p\delta$.

Fix $C=\widehat{\overline{\Q}_p}$, with tilted base $C^\flat$ and
untilt divisor $\infty$.  Let $\mathcal M$ classify $P_3$-structured
forms $E$ of $V_3$ equipped with a quotient line at $\infty$.
The lower modification $E^-$ is a slope-zero rank-three bundle,
hence a rational local system; its determinant inherits an integral
lattice from the $P_3$-structure.  The two frame descriptions give
\[
 \begin{aligned}
 \mathcal M&\simeq[\mathbf P^2_C/P_3]\simeq[\Omega_C^2/J_3],\\
 J_3&=\{g\in\GL_3(\Q_p):v_p(\det g)=0\},
 \end{aligned}
\]
where $\Omega_C^2$ is the projective plane minus its
$\Q_p$-rational hyperplanes.  The first presentation records quotient
lines of a framed $E$; the second records upper modifications of a
rationally framed $E^-$ with the prescribed determinant lattice.
These identifications, including their common arithmetic action, are
proved in \Cref{h3:rat:modification-presentations}.

The projective presentation singles out the reduced-trace line from
$P_3$ as the arithmetic-invariant part of rational degree-three cohomology.
The Drinfeld presentation places a nonzero matrix-trace class in that
same line.  Pulling their proportionality back along the ordinary
sequence then restricts the rank-three matrix trace to the trace of
its original integral rank-two kernel.  We retain these actual
classifying maps through finite Witt coefficients and exact Morava
descent to identify the nonzero image in $\Q_p\delta$ under localization.

\subsection{Integral coefficients and continuous geometric descent}

Put $\Lambda_\ell=\Z/p^\ell$.  For a small v-stack $Z$ in this section, set
\[
 R\Gamma_{\mathrm{int}}(Z)
   =R\varprojlim_{\ell}R\Gamma_v(Z,\Lambda_\ell),\qquad
 H^i_{\mathrm{int}}(Z)=H^i(R\Gamma_{\mathrm{int}}(Z)),
 \qquad H^i_{\mathrm b}(Z)=H^i_{\mathrm{int}}(Z)[1/p].
\]
For a locally profinite group $G$ we use a separate complex in $D(\Z)$,
defined independently of a geometric base:
\[
 R\Gamma_{\mathrm{int}}(G)
    =R\varprojlim_\ell C^\bullet_{\hthreects}(G,\Lambda_\ell),
 \qquad H^i_{\mathrm{int}}(G)=H^i(R\Gamma_{\mathrm{int}}(G)),
 \qquad H^i_{\mathrm b}(G)=H^i_{\mathrm{int}}(G)[1/p].
\]
The coefficients here are finite discrete modules with trivial action,
and the cochains are the inhomogeneous cochains of
\Cref{h3:sec:framework}.  At each finite level the cochain complex is the
external derived cohomology of the condensed classifying topos
$\mathrm B_{\mathrm{cond}}G$, the topos of condensed sets with $G$-action.
Indeed, the standard resolution there has terms indexed by $G^a$.
Finite discrete coefficients are solid, and their higher condensed
cohomology on a locally profinite set vanishes.  Its degree-zero
sections are $C_{\hthreects}(G^a,\Lambda_\ell)$, so the standard
resolution gives the natural identification in the derived category
\[
 R\Gamma_{\mathrm{Cond}}
   (\mathrm B_{\mathrm{cond}}G,\underline{\Lambda_\ell})
       \simeq C^\bullet_{\hthreects}(G,\Lambda_\ell)
\]
by \cite[Section~2, standard resolution and Lemmas~2.1--2.2,
pp.~5--7]{AnschuetzSolidGroupCohomology}.  This is an external complex
of ordinary modules.  We reserve $R\Gamma(G,K)$ for internal condensed
derived invariants of a condensed coefficient object $K$ with its
actual $G$-action; evaluation at the point, denoted $(*)$, gives its
external complex.

Twists are made at finite level.  For compact $G$, reduction of finite
cochains is surjective: lift the finitely many values on their clopen
fibers.  Thus the derived coefficient limit is represented by
$C^\bullet_{\hthreects}(G,\Z_p)$.  Every continuous function from the
compact space $G^a$ to $\Q_p$ is bounded, so exact inversion of $p$
identifies this complex with $C^\bullet_{\hthreects}(G,\Q_p)$.
For the noncompact groups below we retain the displayed definition
using finite coefficients before the derived limit and inversion of $p$.

We now specify how a group class enters relative geometric cohomology.
For a small v-base $B$, let $\underline G_B=B\times\underline G$ be
the constant locally profinite v-group and put
$\mathrm B_B G=[B/\underline G_B]$, with trivial action on $B$.
Products with group parameters below are over this base.

\begin{lemma}[Finite classifying maps]\label{h3:rat:finite-classifying}
For a $G$-space $X$ over $B$ there is a natural finite coefficient map
\begin{equation}\label{h3:eq:rat-finite-classifying}
 \mathsf u_{X,G,\ell}:
 C^\bullet_{\hthreects}(G,\Lambda_\ell)
 \longrightarrow
 \operatorname{Tot}_{[a]\in\Delta}
 R\Gamma_v(X\times\underline{G^a},\Lambda_\ell)
 \simeq R\Gamma_v([X/G],\Lambda_\ell).
\end{equation}
For a $G$-torsor $T\to Z$ over $B$ with relative classifying map
$c_T:Z\to\mathrm B_B G$, this gives
\[
 \operatorname{cl}_{c_T,\ell}
   :=c_T^*\mathsf u_{B,G,\ell}:
 C^\bullet_{\hthreects}(G,\Lambda_\ell)
       \longrightarrow R\Gamma_v(Z,\Lambda_\ell).
\]
It equals \eqref{h3:eq:rat-finite-classifying} for $X=T$ under
$[T/G]\simeq Z$.  These maps commute with continuous group
homomorphisms and extension of torsor structure, equivariant maps,
base change, base automorphisms with compatible descent, coefficient
reductions, and cup products.  They also commute with the finite
coefficient units $\Lambda_\ell=W_\ell(\F_p)\to W_\ell\mathcal O^\flat$.
Their derived coefficient limits and their rationalizations define
maps denoted $\operatorname{cl}_{c_T,\mathrm{int}}$ and
$\operatorname{cl}_{c_T,\mathrm b}$ on the corresponding cohomology.

For an algebraically closed perfectoid field $C$ of characteristic zero, put
$B_C=\operatorname{Spa}(C,\mathcal O_C)^\diamond$.  For every
profinite $Q$ and every finite constant coefficient module $A$ on
$B_C$, the actual constants map is an equivalence
\begin{equation}\label{h3:eq:rat-geometric-point-constants}
 C_{\hthreects}(Q,A)[0]
       \xrightarrow{\ \simeq\ }
 R\Gamma_v(B_C\times\underline Q,A).
\end{equation}
This is natural in $Q$, in $A$, and in base automorphisms with their
actions on $A$.  In particular it includes the finite Tate fibers
$\Lambda_\ell(-j)$.  For compact $G$, it makes
$\mathsf u_{B_C,G,\ell}$ an equivalence.
\end{lemma}

\begin{proof}
For a locally profinite parameter $Q$, regard
$A_Q=C_{\hthreects}(Q,\Lambda_\ell)$ as a discrete ring.  Evaluation
gives a map of ring sheaves on $X\times\underline Q$,
\[
 \underline{A_Q}\longrightarrow\underline{\Lambda_\ell},
 \qquad (f,q)\longmapsto f(q).
\]
It is defined on the finitely many clopen fibers of each $f$; this
uses no compactness of $Q$.  The unit of the derived
constant-sheaf/global-sections adjunction, followed by derived
sections of evaluation, gives
\[
 C_{\hthreects}(Q,\Lambda_\ell)[0]
       \longrightarrow R\Gamma_v(X\times\underline Q,\Lambda_\ell).
\]
Equivalently, this sends a continuous $\Lambda_\ell$-valued function
on $Q$ to its ordinary section and then to derived sections.  With $Q=G^a$, the
projections of the action-groupoid nerve commute with all bar faces
and codegeneracies, including the action face.  These maps therefore
form a cosimplicial map.  Its totalization is
\eqref{h3:eq:rat-finite-classifying} by quotient \v{C}ech descent.

The map of nerves that forgets $T$ over $B$ proves the stated equality
with $c_T^*\mathsf u_{B,G,\ell}$.  For a continuous homomorphism
$h:H\to G$, let $S\to Z$ be an $H$-torsor over $B$ with a specified
identification $S\times^H G\simeq T$ over $Z$.  The resulting map of
torsor nerves acts in degree $a$ by $h^a$ on the transition elements.  Evaluation commutes with these maps, with equivariant
maps of $X$, and with every base change.  A base automorphism whose
descent commutes with the constant group acts on a torsor nerve by
$(t,g_1,\ldots,g_a)\mapsto(\gamma(t),g_1,\ldots,g_a)$, which proves
the asserted arithmetic naturality.  The construction uses maps of
coefficient ring sheaves and the lax monoidal derived sections functor,
so it preserves cup products.  Composing the same ring-sheaf maps
with $\Lambda_\ell\to\Lambda_{\ell-1}$ or with
$\Lambda_\ell=W_\ell(\F_p)\to W_\ell\mathcal O^\flat$ proves the
coefficient naturalities term by term.  In particular these are maps
of the derived coefficient towers before cohomology is taken.

For \eqref{h3:eq:rat-geometric-point-constants}, use all finite clopen
quotients $Q_i$ of $Q$.  The example after Definition~7.8 and
Proposition~7.16 of \cite{ScholzeDiamonds} identify
$B_C\times\underline Q$ with the inverse limit of $B_C\times Q_i$
and make $B_C$ a strictly totally disconnected geometric point.
Its \'etale covers split, so its finite constant cohomology is $A[0]$.
Continuity and comparison with v-cohomology in
\cite[Propositions~14.9--14.10]{ScholzeDiamonds} now give
\[
 R\Gamma_v(B_C\times\underline Q,A)
   \simeq\colim_i\hthreeMap(Q_i,A)[0]
   =C_{\hthreects}(Q,A)[0].
\]
Every map here is induced by the coefficient unit and is natural in
the finite partitions, coefficient fiber, and base.  This proves the
stated naturality, including the arithmetic action on a Tate fiber.
Apply the equivalence to $Q=G^a$ for compact $G$ and totalize to obtain
the last assertion.  Over the finite-field base used below we use the
finite classifying map itself; this point comparison is only asserted
over $B_C$.
\end{proof}

We retain profinite parameters throughout geometric descent.  Write
\[
 \mathscr C_{Z,\ell}(Q)
   =R\Gamma_v(Z\times\underline Q,\Lambda_\ell),\qquad
 \mathscr C_Z=R\varprojlim_\ell\mathscr C_{Z,\ell},
\]
where $Q$ is profinite; these are derived sections on $Q$.

\subsection{Trace classes and a parabolic subgroup}

The invariant form
\[
 (X,Y,Z)\longmapsto\hthreeTr(X[Y,Z])
\]
uses reduced trace for the division algebra and matrix trace for the
split groups.  This is the classical degree-three cocycle attached to
an invariant symmetric form, as in Chevalley--Eilenberg
\cite[Section~21]{ChevalleyEilenberg1948} and Koszul
\cite[Section~11]{Koszul1950}.
The next lemma constructs the corresponding group-cohomology classes
and identifies the restriction needed for the ordinary kernel.

\begin{lemma}\label{h3:rat:trace-groups}
The groups $P_3$, $J_3$, and $\GL_2(\Z_p)$ have rational cohomology
$H^1_{\mathrm b}=\Q_p$, $H^2_{\mathrm b}=0$, and
$H^3_{\mathrm b}=\Q_p$.  They admit nonzero classes
$e_3$, $e_J$, and $e_2$, respectively, whose rational Lie-algebra
images are nonzero multiples of the displayed trace form.
For every compact open $K\subset J_3$, restriction followed by
the compact Lie-algebra comparison identifies $H^*_{\mathrm b}(J_3)$
with $H^*(\mathfrak{gl}_3,\Q_p)$; the identifications for different
$K$ agree under further restriction.  For
\[
 D=\left\{\begin{pmatrix}A&v\\0&a\end{pmatrix}:
 A\in\GL_2(\Z_p),\ a\in\Z_p^\times,\ v\in\Q_p^2\right\},
\]
write $\iota:D\to J_3$ and $\rho:D\to\GL_2(\Z_p)$ for
inclusion and projection.  Then
\begin{equation}\label{h3:eq:rat-parabolic-trace}
 \iota^*e_J=b\rho^*e_2\quad\text{for some }b\in\Q_p^\times.
\end{equation}
\end{lemma}

\begin{proof}
We first make the continuous comparison on compact groups precise.
Each compact group $K$ under consideration has an open normal uniform
pro-$p$ subgroup $U$.  For $P_3$ and $\GL_2(\Z_p)$ one may use a
sufficiently deep principal congruence subgroup.  A compact open
$K\subset J_3$ preserves a lattice in $\Q_p^3$: the orbit of any
lattice is finite, and its sum is stable.  A sufficiently deep
principal congruence subgroup for this lattice lies in $K$ and is
normal there.  If the depth is at least one, the corresponding Lie
lattice $\mathfrak u=p^m\mathcal O$, for the relevant endomorphism
order $\mathcal O$, satisfies
$[\mathfrak u,\mathfrak u]\subset p\mathfrak u$.  At odd $p$,
logarithm and exponential converge on these lattices, identify the
$p$th-power map with multiplication by $p$, and show that the group
is finitely generated, powerful, and torsion free.  It is therefore
uniform.

Give $U$ its lower-$p$-series valuation.  For $p\geq5$ it is saturated
and equi-$p$-valued, with value denominator $e=1$ and graded generators
in degree one.  The trivial finite free coefficient module $\Z_p$
has action in $1+p\operatorname{End}_{\Z_p}(\Z_p)$.  The continuous
comparison of \cite[Theorem~3.3.3]{HuberKingsNaumann} therefore gives
a cup-compatible isomorphism
\[
 H^*_{\hthreects}(U,\Z_p)\simeq H^*(\mathfrak u,\Z_p),
\]
where $\mathfrak u$ is the integral Lie algebra of $U$.
Coefficient naturality and rational agreement in
\cite[Theorem~3.1.1(1)--(2)]{HuberKingsNaumann}, together with the
boundedness of compact rational cochains, identify its rationalization
with Lazard's comparison for $\operatorname{Lie}(K)=\mathfrak u[1/p]$.
Continuous group homomorphisms preserve lower $p$-series, so
Theorem~3.1.1(3) of the same source gives naturality for restrictions
and for the block maps after passing to such small opens.

The rational Hochschild--Serre sequence for $U\triangleleft K$
collapses to finite quotient invariants: positive cohomology of the
finite group $K/U$ vanishes by averaging in $\Q_p$.  Here
$\operatorname{Lie}(K)$ is the Lie algebra of a $\GL_n$-form, with $n=2$ or $3$.
The exterior calculation and triviality of its adjoint action are
given by \cite[Lemma~3.8.2]{BSSW}.  They give generators in degrees
$1,3,\ldots,2n-1$, without choosing explicit cocycles for all of them.
Thus the finite quotient fixes the cohomology, and in particular the
low-degree dimensions are $1,0,1$ in degrees $1,2,3$.

We identify the particular degree-three generator needed here.
For the split Lie algebra put $q(X,Y)=\operatorname{Tr}(XY)$;
for the inner form use reduced trace.  Cyclicity gives
$q([X,Y],Z)=\operatorname{Tr}(X[Y,Z])$.  After an algebraically
closed splitting-field extension the restriction of $q$ to
$\mathfrak{sl}_n$, for $n=2,3$, is nonzero and nondegenerate.
The degree-three invariant-form calculation of
\cite[Section~11, Theorems~11.1--11.2]{Koszul1950}
therefore makes its class nonzero.  The finite Chevalley--Eilenberg
cochain complex commutes with field extension, so the class already
defined on the form over $\Q_p$ is nonzero.  Comparing restrictions through a
common small uniform subgroup shows that the group identifications
are independent of the chosen open subgroup.  The upper-left block
map restricts the displayed degree-three matrix trace form to the
identical form on the $2\times2$ block, and hence restricts it
nontrivially.

For $P_3$, choose $e_3$ to have exactly this reduced-trace Lie-algebra
class under the natural compact comparison.  The coefficient-field
group $\Gamma=\Gal(k/\F_p)$ acts by automorphisms of the Honda division
algebra and preserves reduced trace.  Naturality of that comparison
therefore makes this normalized rational class $\Gamma$-invariant.

For the noncompact group $J_3$, use the reduced affine building in
\Cref{h3:lem:rank-three-building} with $F=\Q_p$.  That lemma supplies
a contractible two-dimensional complex on which $J_3$ preserves types,
has compact open face stabilizers fixing their faces pointwise, and has
a closed chamber as a fundamental domain.  The cellular resolution and continuous
Shapiro give, first with finite coefficients and then with the
derived integral limit,
\[
 E_1^{a,s}=\bigoplus_{\substack{\sigma\subset\Delta\cr
                         \dim\sigma=a}}
 H^s_{\mathrm{int}}(K_\sigma)
       \ \Longrightarrow\ H^{a+s}_{\mathrm{int}}(J_3).
\]
There are finitely many summands and cellular degrees.  After
inversion of $p$, every restriction in a fixed row identifies with
the identity on $H^s(\mathfrak{gl}_3,\Q_p)$; changing a face
representative only introduces an inner conjugation.  That row
is therefore the cellular cochain complex of the simplex $\Delta$
with these constant coefficients.  It has cohomology only in
cellular degree zero.  The edge map, and then further restriction
to an arbitrary compact open subgroup using intersections, gives
the required assertion for $J_3$.

It remains to account for the unipotent radical of $D$.
Fix $\ell$ and write $\Q_p^2=\bigcup_{r\geq0}p^{-r}\Z_p^2$.
The two-variable Koszul resolution computes
\[
 H^q(p^{-r}\Z_p^2,\Z/p^\ell)
       =\bigwedge^q(\Z/p^\ell)^2.
\]
Restriction from stage $r+1$ to $r$ multiplies degree $q$ by
$p^q$.  These systems are pro-zero for $q>0$ and constant for
$q=0$.  The restrictions of cochains themselves are surjective:
the smaller subgroup is clopen, so a cochain extends by zero.
Thus their inverse limit computes the derived limit, and the
Milnor sequence gives
$R\Gamma(\Q_p^2,\Z/p^\ell)=\Z/p^\ell[0]$.
The calculation remains valid with profinite parameters, by the
same finite Koszul complexes and extension of cochains.
Thus this equality holds for internal derived condensed invariants.
The constants map is equivariant for the Levi action, and its
underlying equivalence is therefore an equivariant equivalence.
Hochschild--Serre consequently makes inflation from the Levi subgroup
an equivalence at each finite level and then integrally.
The rational K\"unneth calculation for
$\GL_2(\Z_p)\times\Z_p^\times$ has degree three equal to the
degree-three group of its first factor, since the second factor
has cohomological dimension one and the first has $H^2=0$.
Restriction along $A\mapsto\operatorname{diag}(A,1)$ detects the
nonzero trace form.  This proves \eqref{h3:eq:rat-parabolic-trace}.
\end{proof}

\subsection{Two presentations of a modification}

Write
\[
 B_k=\Spd k,\qquad
 B_C=\Spd C^\flat\simeq\operatorname{Spa}(C,\mathcal O_C)^\diamond.
\]
The chosen embedding $F=W(k)[1/p]\subset C$ gives the embedding
$k\subset C^\flat$ and hence the base map $B_C\to B_k$.
All quotient presentations in this subsection are over $B_C$.
A determinant lattice on a slope-zero rank-three bundle is a
$\Z_p$-lattice in its determinant rational local system.  Its frame
group is $J_3$.  Fix a lattice in the one-dimensional rational
local system associated to
\[
 \Delta_3=\det(V_3)(-\infty).
\]
The action of $D_3^\times$ on it is reduced norm.  Its stabilizer
is therefore exactly $P_3$.

\begin{lemma}\label{h3:rat:modification-presentations}
Let $\mathcal M$ classify a form of $V_3$ with $P_3$-structure
and a quotient line at $\infty$.  There are equivalences
\begin{equation}\label{h3:eq:rat-two-presentations}
 \mathcal M\simeq[\mathbf P^2_C/P_3]
                 \simeq[\Omega_C^2/J_3].
\end{equation}
The relative classifying maps
\[
 c_{3,\mathcal M}:\mathcal M\longrightarrow\mathrm B_{B_C}P_3,
 \qquad
 c_{J,\mathcal M}:\mathcal M\longrightarrow\mathrm B_{B_C}J_3
\]
record the positive bundle and its lower modification, respectively.
For the prescribed finite unramified field $F=W(k)[1/p]$, these
presentations carry a common semilinear arithmetic action of $G_F$
over its action on $B_C$, commuting with both constant groups and
inducing the usual action on $\Omega_C^2$.
\end{lemma}

\begin{proof}
The kernel of a length-one quotient of a form $E$ of $V_3$ is
a rank-three bundle $E^-$ of degree zero.  If it had a positive
Harder--Narasimhan subbundle of rank $b<3$ and integral degree
$d\geq1$, its inclusion in $E$ would contradict semistability:
$d/b\geq1/b>1/3$.  Thus $E^-$ has slope zero.  By the relative
slope-zero correspondence it is a rational local system
\cite[Corollary~8.7.10]{KedlayaLiu}; see \Cref{h3:geom:curve-inputs}.

Conversely an upper length-one modification of $\mathcal O^3$
has, at a geometric point with local parameter $t$ at $\infty$,
lattice
\[
 (B_{\mathrm{dR}}^+)^3+
 B_{\mathrm{dR}}^+t^{-1}\widetilde\ell
\]
for a direction line $\ell$.  If $\ell$ lies in a proper
rational subspace $W$, modifying $W\otimes\mathcal O$ gives a
degree-one subbundle of rank less than three, which destabilizes.
In the other direction, let $F'$ be a saturated destabilizing
subbundle of the upper modification, of rank $b<3$ and degree
$d\geq1$.  Its intersection $F''$ with $\mathcal O^3$ has
degree $d-\epsilon$, where $0\leq\epsilon\leq1$.  Semistability
of $\mathcal O^3$ forces $d=\epsilon=1$ and $\deg F''=0$.
The saturation of $F''$ in $\mathcal O^3$ has degree at least zero,
while semistability forces its degree to be at most zero.  Its torsion
quotient therefore has degree zero, so $F''$ is already saturated.
Thus $F''$ is a
slope-zero subbundle of $\mathcal O^3$, hence equals
$W\otimes\mathcal O$ for a rational subspace $W$; maps between
trivial bundles are matrices over $H^0(\mathcal O)=\Q_p$.
The equality $\epsilon=1$ says that $\ell\subset W_C$.
This proves precisely the Drinfeld condition.

Let $\mathcal Z$ be the sheaf of injections
$\mathcal O^3\hookrightarrow V_3$ with length-one cokernel at
$\infty$ preserving the chosen determinant lattices.  The
commuting actions are postcomposition by $P_3$ and inverse
precomposition by $J_3$.  The preceding slope calculation and
the relative trivial/basic-stratum theorems
\cite[Theorems~III.2.4 and~III.4.5]{FarguesScholze} give actual sheaves of
frames in families.  They identify
$\mathcal Z/J_3=\mathbf P^2_C$.  In the other direction they
identify $\mathcal Z/P_3=\Omega_C^2$: the full frame group of
the positive bundle is $D_3^\times$, and its surjective norm
valuation lets one match the determinant lattice locally.
Taking the remaining quotient proves
\eqref{h3:eq:rat-two-presentations}.

The arithmetic field is $F=W(k)[1/p]$; the coefficient field
of the curve remains $\Q_p$.  Let $N$ be the fixed Honda
isocrystal over $k$, whose period-ring construction gives
$V_3$.  All Honda endomorphisms are defined over $k$.
On the period coefficients tensored over $F$ with $N$, define
the action of $\gamma\in G_F$ by $\gamma\otimes\hthreeid_N$.
It commutes with Frobenius and so descends to the bundle.
Every element of $D_3$ is an $F$-linear endomorphism of $N$
commuting with its Frobenius; its matrix entries are fixed
by $G_F$.  Thus this specified descent commutes with the
entire $D_3$-action.  We have not trivialized a torsor of
frames of an arbitrary descended basic bundle.

The untilt ideal $\ker\theta$ is functorial, so $\Delta_3$
inherits descent.  The relative slope-zero correspondence
makes its rational section space a continuous one-dimensional
representation of $G_F$.  This representation need not be
trivial.  Its valuation character has compact subgroup image
in $\Z$, hence zero, so every lattice in it is stable;
no generator is asserted to be invariant.
Define the arithmetic action on an injection $u$ by pullback
and the specified descent on source and target.  The standard
rational coordinate vectors of $\mathcal O^3$ are fixed, so
\[
 \gamma(guj^{-1})=g\gamma(u)j^{-1}
       \qquad(g\in P_3,\ j\in J_3).
\]
The determinant condition is preserved by lattice stability.
On the upper side, directions lie in
$\Q_p^3\otimes(\mathcal O(\infty)/\mathcal O)|_\infty$.
The semilinear scalar in the second factor is common to all
coordinates and cancels in projective space.  This is the usual
arithmetic action on $\Omega_C^2$, commuting with $J_3$.
No arithmetic invariant choice of a frame is required.
\end{proof}

\begin{lemma}\label{h3:rat:fixed-projective}
The first presentation in \eqref{h3:eq:rat-two-presentations} has
a natural arithmetic-equivariant splitting
\begin{equation}\label{h3:eq:rat-projective-h3}
 H^3_{\mathrm b}(\mathcal M)
   =H^3_{\mathrm b}(P_3)\oplus H^1_{\mathrm b}(P_3)(-1),
\end{equation}
whose first inclusion is the classifying map
$\operatorname{cl}_{c_{3,\mathcal M},\mathrm b}$ of
\Cref{h3:rat:finite-classifying}.
\end{lemma}

\begin{proof}
For a profinite set
$Q=\varprojlim_iQ_i$ with finite $Q_i$, continuity on spatial
diamonds and comparison of finite coefficients give
\[
 R\Gamma_v(\mathbf P^2_C\times\underline Q,\Lambda_\ell)
  \simeq\colim_i\hthreeMap\bigl(Q_i,
                 R\Gamma_{\acute et}(\mathbf P^2_C,\Lambda_\ell)\bigr).
\]
Here one applies \cite[Proposition~14.9]{ScholzeDiamonds} to
$\mathbf P^2_C\times Q_i$ and then
\cite[Proposition~14.10 and Lemma~15.6]{ScholzeDiamonds}; these statements
allow the coefficient ring $\Lambda_\ell$ also when its torsion
prime is $p$.  The finite-coefficient cohomology of the analytic
projective space agrees with its algebraic cohomology: the
proper comparison of \cite[Theorem~3.12]{ScholzeSurvey},
restating \cite[Theorem~3.7.2]{Huber}, applies to
$\mathbf P^2_C\to\Spec C$ with $\F_p$ coefficients, and
the coefficient exact sequences extend it to $\Lambda_\ell$.
Thus the projective-space calculation is valid as the actual
condensed geometric complex.

The anticanonical bundle has canonical projective equivariance,
including arithmetic descent, and restricts to $\mathcal O(3)$.
Since $p\geq5$, the class
$h_\ell=3^{-1}c_1(\omega_{\mathbf P^2}^{-1})$
on the quotient restricts to the hyperplane class.  Its powers
give an actual morphism of equivariant derived coefficient objects
\[
 \bigoplus_{j=0}^2\Lambda_\ell(-j)[-2j]
       \longrightarrow\mathscr C_{\mathbf P^2_C,\ell}.
\]
The preceding calculation shows that this is an equivalence on
every profinite parameter.  Apply it to $Q=P_3^a$ in the quotient
nerve and totalize.  The finite direct sum commutes with
totalization, giving an equivalence
\[
 \bigoplus_{j=0}^2
 C^\bullet_{\hthreects}(P_3,\Lambda_\ell(-j))[-2j]
       \xrightarrow{\ \simeq\ }
 R\Gamma_v(\mathcal M,\Lambda_\ell).
\]
In its $j=0$ component, each bar term sends a continuous constant
function to its coefficient-unit section.  By
\Cref{h3:rat:finite-classifying}, this component is exactly
$\operatorname{cl}_{c_{3,\mathcal M},\ell}$.
The Kummer classes commute with coefficient reduction.  Take the
derived limit in $\ell$ and then invert $p$; the finite direct sum
commutes with both operations.  Only $j=0,1$ contribute to degree
three, giving \eqref{h3:eq:rat-projective-h3}.  Naturality of the
Kummer classes gives its arithmetic compatibility and the
displayed twists.
\end{proof}

The projective splitting identifies the classifying image of the
$P_3$ trace.  We next need to show that the $J_3$ trace has a
nonzero classifying image on the same modification stack.  The
Drinfeld cohomology theorem gives arithmetic scalar actions on
point-valued cohomology; the following finite resolution justifies
their use on the compact-group descent rows.

\begin{lemma}[Finite resolutions for geometric coefficients]
\label{h3:rat:solid-resolution}
The complexes $\mathscr C_Z$ are bounded below, derived solid, and
derived $p$-complete.  Let $G$ be a compact $p$-adic analytic group
without $p$-torsion.  There is a finite projective resolution
$P_\bullet\to\Z_p$ over $\Lambda_G=\Z_p[[G]]$, with finitely
generated terms and length $\dim G$, such that, for every derived
solid $\underline{\Z_p}$-module complex $K$ with
$\underline{\Z_p}$-linear continuous $G$-action,
\begin{equation}\label{h3:eq:rat-solid-hom}
 R\Gamma(G,K)(*)\simeq\Hom_{\Lambda_G}(P_\bullet,K).
\end{equation}
Here $\Hom$ is the external complex of morphisms in solid
$\underline{\Lambda_G}$-modules.  A continuous action means an
actual module over $\underline{\Z_p}[\underline G]$ in the
derived condensed category, rather than an action on $K(*)$.
If $K$ is bounded below, and an operator $\gamma$ commuting with $G$
acts by the scalar $\lambda_j\in\Q_p$ on $H^j(K)(*)[1/p]$, it
acts by $\lambda_j$ on
$H^s(G,H^j(K))(*)[1/p]$ for every $s$.  For $K=\mathscr C_Z$ the
resulting spectral sequence is
\begin{equation}\label{h3:eq:rat-solid-descent}
 H^s(G,H^j(\mathscr C_Z))(*)[1/p]
       \ \Longrightarrow\ H^{s+j}_{\mathrm b}([Z/G]),
\end{equation}
with a finite filtration in each total degree.
\end{lemma}

\begin{proof}
Let $V=\operatorname{Spa}(R,R^+)$ be affinoid perfectoid, and write
$R^\flat$ for its tilt (equal to $R$ in characteristic $p$).
Its product with $\underline Q$ is affinoid perfectoid with ring
$C_{\hthreects}(Q,R)$ and tilt $C_{\hthreects}(Q,R^\flat)$.
Affinoid acyclicity gives
\[
 R\Gamma_v(V\times\underline Q,\mathcal O^\flat)
       =C_{\hthreects}(Q,R^\flat)[0]
\]
by \cite[Proposition~8.8]{ScholzeDiamonds}.  Choose a decreasing countable basis of
open additive subgroups $R^\flat_a$ of the complete additive group $R^\flat$.
Then $\underline{R^\flat}=\varprojlim_a(R^\flat/R^\flat_a)_{\mathrm{disc}}$ is solid.
This limit is also the derived limit: continuous maps from a
profinite set to each discrete quotient have finite image, and
lifts of their finitely many values give surjectivity at each
successive quotient.  A v-hypercover by disjoint unions of affinoid
perfectoids now expresses the $\mathcal O^\flat$ complex, with all
parameters, as a limit of products of solid modules.  Derived solid
modules are closed under limits and extensions
\cite[Theorem~4.4(ii)]{Tang}.  Artin--Schreier gives the assertion
for $\F_p$, the coefficient exact sequences give it for $\Z/p^\ell$,
and the derived limit gives it for $\mathscr C_Z$.  These complexes
are bounded below.  The finite coefficient complexes are derived
$p$-complete, so their limit is too.  Evaluation at the point is
exact and preserves limits; in particular
\[
 H^j(\mathscr C_Z)(*)=H^j(R\Gamma_{\mathrm{int}}(Z)).
\]

The compact-resolution argument in \Cref{h3:co:compact-resolution} gives
Noetherianity of $\Lambda_G$ and projective dimension $\dim G$ for
its trivial profinite module, using
\cite[Sections~1.1--1.2, pp.~275--276]{Venjakob} together with
\cite[Section~1, Corollary~(1)]{Serre1965CohomologicalDimension}.
Successive finite free presentations therefore give the stated
finite resolution.  We give the derived comparison needed to use
it.  Let $\mathcal C=D(\operatorname{Cond}_{\underline{\Z_p}})$
and $\mathcal D=D(\operatorname{Solid}_{\underline{\Z_p}})$.
Solidification is a symmetric monoidal localization
$L:\mathcal C\to\mathcal D$ with fully faithful right adjoint
$I$ \cite[Theorem~4.4(ii)]{Tang}.
For the free condensed group ring $A=\underline{\Z_p}[\underline G]$,
its solidification $LA$ is $\underline{\Lambda_G}$ in degree zero,
including multiplication \cite[Section~3.6 and the proof of
Theorem~4.4]{Tang}.
The underived assertion also holds for a free module on any
profinite set $Q$: choose an extremally disconnected hypercover
$Q_\bullet\to Q$ and solidify its free resolution.  The resulting
augmented complex is $\Z_p[[Q_\bullet]]\to\Z_p[[Q]]$.
Its Pontryagin dual is the augmented complex of continuous
$\Q_p/\Z_p$-valued functions.  Discrete coefficients are acyclic
on a profinite set, by passage to finite clopen partitions, so
that dual complex is exact.  Pontryagin duality and exact
condensation prove the assertion.  Functoriality with respect
to multiplication on $G$ gives the ring identification.

For any symmetric monoidal localization there is an induced
adjunction
\[
 L_A:\hthreeMod_A(\mathcal C)\ \rightleftarrows\
       \hthreeMod_{LA}(\mathcal D):J,
\]
where $J$ is inclusion followed by restriction along
$A\to I(LA)$.  Its unit has underlying map $M\to ILM$,
and its counit has underlying map $LIN\to N$.
The counit is an equivalence; hence $J$ is fully faithful,
with essential image exactly the modules whose underlying
complex is derived solid.  This proves the assertion in
unbounded degrees.  To identify the target, apply this same
localization to the already solid ring
$B=\underline{\Lambda_G}$.  Its essential image inside
$D(\operatorname{Cond}_B)$ consists of precisely the
complexes with derived solid underlying coefficients.
By \cite[Theorem~4.4]{Tang}, the unbounded category
$D(\operatorname{Solid}_B)$ has that same essential image.
This identifies the two derived module categories; it is
stronger than an equivalence of their hearts.  Equivalently,
the comparison is computed on free modules $B[Q]$ for
extremally disconnected profinite $Q$ and their solidifications,
which give the projective resolutions in these categories.
Thus the actual condensed action on $K$ belongs to this
category, and full faithfulness identifies its external
derived invariants with those computed there.  For geometric
$K$, the action and its coherent group law come from the
maps on $Z\times\underline Q$ with all group parameters
retained; the preceding solidity argument applies to its
underlying complex.

Condensation preserves the exact profinite resolution
\cite[Theorem~3.2]{Tang}.  Every $P_a$ is a retract of a finite
free $\Lambda_G$-module, and
$\Hom_{\Lambda_G}(\Lambda_G,M)=M(*)$ is exact.  Thus these terms
are projective against arbitrary solid coefficient modules and
prove \eqref{h3:eq:rat-solid-hom}.  The complex
$\Hom_{\Lambda_G}(P_\bullet,H^j(K))$ consists of retracts of
finite sums of $H^j(K)(*)$.  Its retracts commute with $\gamma$.
After exact inversion of $p$, the entire complex therefore has
scalar action $\lambda_j$, proving the assertion about its
cohomology.  No assertion about a condensed sheaf is inferred from
its point value.

Finally the nerve of $Z\to[Z/G]$ has terms
$Z\times\underline G^a$.  With parameters retained, its
descent complex is precisely the continuous bar construction on
$\mathscr C_Z$.  At finite level the equivariant coefficient unit
$\underline{\Lambda_\ell}\to\mathscr C_{Z,\ell}$ is the evaluation
map of \Cref{h3:rat:finite-classifying} on every profinite parameter.
Its map on $G$-invariants is therefore exactly
$\mathsf u_{Z,G,\ell}$, since their degree-$a$ bar maps agree.
This identity also holds after the derived coefficient limit and
inversion of $p$.  Its source after that limit is
$R\Gamma(G,\underline{\Z_p})(*)\simeq R\Gamma_{\mathrm{int}}(G)$:
the finite reductions of constants are surjective on every profinite
parameter, and internal invariants and point evaluation preserve
limits.  The resolution bounds the group direction by
$\dim G$, giving \eqref{h3:eq:rat-solid-descent} and the asserted
finite filtration before and after rationalization.
\end{proof}

\begin{lemma}\label{h3:rat:two-traces}
On $\mathcal M$, abbreviate the classifying images
$\operatorname{cl}_{c_{3,\mathcal M},\mathrm b}(e_3)$ and
$\operatorname{cl}_{c_{J,\mathcal M},\mathrm b}(e_J)$ by $e_3$ and
$e_J$.  Both are nonzero and satisfy $e_3=\kappa e_J$ for some
$\kappa\in\Q_p^\times$.  We retain this abbreviation for subsequent
pullbacks from $\mathcal M$.
\end{lemma}

\begin{proof}
By \eqref{h3:eq:rat-projective-h3}, the class $e_3$ is nonzero.
Arithmetic acts trivially on its first summand and by
$\chi_{\mathrm{cyc}}^{-1}$ on the second.  The cyclotomic
character of $G_F$ has open image, so the invariant subspace
is precisely $\Q_p e_3$.

Choose a compact uniform open subgroup $U\subset J_3$.  The
restriction $\operatorname{res}_U(e_J)$ of the source group class
$e_J\in H^3_{\mathrm b}(J_3)$ is nonzero by \Cref{h3:rat:trace-groups}.
Let $q_U:[\Omega_C^2/U]\to\mathcal M$ be the quotient map and
$c_U:[\Omega_C^2/U]\to\mathrm B_{B_C}U$ its relative classifying map.
Extension of torsor structure gives
$c_{J,\mathcal M}q_U\simeq(\mathrm B_{B_C}(U\hookrightarrow J_3))c_U$.
Thus \Cref{h3:rat:finite-classifying} identifies the restriction of
the geometric $e_J$ with
$\operatorname{cl}_{c_U,\mathrm b}(\operatorname{res}_U(e_J))$.
Apply \eqref{h3:eq:rat-solid-descent} to $[\Omega_C^2/U]$.
The integral Drinfeld computation
\cite[Theorems~1.1 and~5.1]{CDN} identifies
$H^j(\mathscr C_{\Omega_C^2})(*)[1/p]$ as a representation
on which arithmetic acts by the scalar
$\chi_{\mathrm{cyc}}^{-j}$, with $j=0,1,2$ and no higher
rows.  Its integral convention is the derived limit of finite
$p$-power coefficients followed by inversion of $p$, exactly
the convention here: the finite reduction tower defines
$\widehat{\Z}_p$ on the pro-etale site, and derived global
sections preserve its derived inverse limit.  The calculation is for hyperplanes
rational over the ground field, here $\Q_p$; we merely restrict
its arithmetic action to $G_F$.

In degree zero the same computation gives a one-dimensional rational
point value.  The finite coefficient units form a map
$\eta:\underline{\Z_p}\to\mathscr C_{\Omega_C^2}$ after their derived
limit: on every profinite parameter the finite reductions of constants
are surjective, so their limit is $\underline{\Z_p}$.
Pullback to any geometric point sends the unit to $1$.  Hence
\[
 H^0(\eta)(*)[1/p]:\Q_p
       \xrightarrow{\ \simeq\ }
 H^0(\mathscr C_{\Omega_C^2})(*)[1/p].
\]
For each finite projective term $P_a$ in the $U$-resolution of
\Cref{h3:rat:solid-resolution}, applying $\Hom_{\Lambda_U}(P_a,-)$
gives a retract of a finite sum of these point maps, and the unit
commutes with its defining idempotent.  It is therefore an isomorphism
after inversion of $p$ in every term.  Naturality of the descent
spectral sequence for $\eta$ maps its single constant row to the
bottom row of geometric descent.  By the compact bar identity, the
$E_2^{3,0}$ bottom-row map of
the spectral-sequence morphism induced by the finite units defining
$\operatorname{cl}_{c_U,\mathrm b}$ is the resulting isomorphism
\[
 H^3_{\mathrm b}(U)
       \xrightarrow{\ \simeq\ }
 H^3(U,H^0(\mathscr C_{\Omega_C^2}))(*)[1/p].
\]
It sends $\operatorname{res}_U(e_J)$ to the
nonzero bottom-row class.  This argument uses the rational point
value and the natural finite projective complex.

By \Cref{h3:rat:solid-resolution}, an element
$\gamma\in G_F$ with cyclotomic character of infinite order
acts on the actual $j$th descent row by this scalar.  The only
possible incoming differentials to $(3,0)$ are from $(1,1)$
and $(0,2)$.  Their source scalars differ from $1$, so both
differentials vanish.  There is no outgoing differential from
the bottom row.  Thus the nonzero bottom-row image of
$\operatorname{res}_U(e_J)$ survives in
$H^3_{\mathrm b}([\Omega_C^2/U])$, proving that the geometric class
$e_J$ on $\mathcal M$ is nonzero.  The map to $\mathrm B_{B_C}J_3$ is
arithmetic equivariant, so this class is invariant and belongs
to the line $\Q_p e_3$.  This proves the assertion.
\end{proof}

\subsection{The actual rank-two kernel on the ordinary stratum}

Let $\mathcal Y_1$ be the ordinary stratum stack of
\Cref{h3:prop:integral-frames}, over $B_k$.  Put
$G_o=\GL_2(\Z_p)\times P_1$.  Its two frame presentations give
relative classifying maps
\[
 c_{3,k}:\mathcal Y_1\longrightarrow\mathrm B_{B_k}P_3,
 \qquad c_{o,k}:\mathcal Y_1\longrightarrow\mathrm B_{B_k}G_o.
\]
On this stack, $e_3$ means $\operatorname{cl}_{c_{3,k},\mathrm b}(e_3)$,
and $e_2$ means
$\operatorname{cl}_{c_{o,k},\mathrm b}(\operatorname{pr}_1^*e_2)$,
where $\operatorname{pr}_1:G_o\to\GL_2(\Z_p)$ is the first projection.
Set $\mathcal Y_{1,C}=\mathcal Y_1\times_{B_k}B_C$ and write
$c_{3,C},c_{o,C}$ for the base-changed classifying maps.  The finite
base-change identity in \Cref{h3:rat:finite-classifying} identifies
the corresponding classes on $\mathcal Y_{1,C}$ with the base changes
of these same classes over $k$.

\begin{proposition}[Transport along the ordinary kernel]
\label{h3:prop:trace-transport}
There is $a\in\Q_p^\times$ such that
\begin{equation}\label{h3:eq:rat-ordinary-traces}
 e_3=a e_2\quad\text{in }H^3_{\mathrm b}(\mathcal Y_{1,C}).
\end{equation}
The class $e_2$ is pulled back from the actual integral frame
group of the rank-two kernel.
\end{proposition}

\begin{proof}
The universal sequence on this stack is
\[
 0\longrightarrow\mathcal L_2\otimes_{\Z_p}\mathcal O
 \longrightarrow\mathcal V_3\longrightarrow\mathcal V_1
 \longrightarrow0,
\]
where $\mathcal L_2$ is the integral rank-two local system of
\Cref{h3:prop:integral-frames}.  The canonical quotient
$\mathcal V_1\to i_{\infty*}i_\infty^*\mathcal V_1$ defines
$\mathcal V_1^- =\mathcal V_1(-\infty)$ and, by pullback,
a lower modification $\mathcal V_3^-$ of $\mathcal V_3$.
It gives a map $m:\mathcal Y_{1,C}\to\mathcal M$ and the exact
sequence of slope-zero bundles
\[
 0\longrightarrow\mathcal L_2\otimes\mathcal O
 \longrightarrow\mathcal V_3^-
 \longrightarrow\mathcal V_1^-\longrightarrow0.
\]
Apply the relative slope-zero equivalence on the pro-etale
site of a perfectoid test object over $\mathcal Y_{1,C}$.
To verify exactness without assuming it, take a common
pro-etale cover trivializing the three rational local systems.
Full faithfulness identifies the maps with matrices of
sections of $\underline{\Q_p}$.  Write the quotient map as
$(b_1,b_2,b_3)$.  The open loci $b_i\ne0$ cover, since its
rank is one on every geometric fiber.  On such a locus,
$1\mapsto b_i^{-1}e_i$ is a continuous section.  The kernel
has basis $e_j-(b_j/b_i)e_i$ for $j\ne i$.  The given map
from $\mathcal L_2\otimes\Q_p$ to this kernel is a square
matrix invertible on every geometric fiber; its inverse is
continuous by the determinant formula.  Thus the sequence
of rational local systems is exact and locally split.

The $P_3$-structure gives an integral determinant lattice
on the middle local system.  Its tensor quotient by
$\det\mathcal L_2$ is an integral lattice on the quotient
line.  Refine the cover to choose an integral basis of the
original $\mathcal L_2$ and an integral generator of this
quotient lattice, and lift the generator by the displayed
splitting.  The resulting middle frame has determinant
generating its prescribed determinant lattice.  Its possible
changes are exactly
$\bigl(\begin{smallmatrix}A&v\\0&a\end{smallmatrix}\bigr)$
with $A\in\GL_2(\Z_p)$, $a\in\Z_p^\times$, and
$v\in\Q_p^2$.  Its determinant is a unit, so these adapted
frames form exactly the $D$-torsor of \Cref{h3:rat:trace-groups}.
Forgetting the quotient generator and its lift gives the
original integral Tate-frame torsor under $D\to\GL_2(\Z_p)$.
These constructions commute with all perfectoid base changes
and descend on the v-site, giving the actual relative map
$c_D:\mathcal Y_{1,C}\to\mathrm B_{B_C}D$.  The identifications of
the associated torsors give homotopies of relative classifying maps
\[
 \begin{aligned}
 c_{3,\mathcal M}m&\simeq c_{3,C},\\
 (\mathrm B_{B_C}\iota)c_D&\simeq c_{J,\mathcal M}m,\\
 (\mathrm B_{B_C}\rho)c_D
     &\simeq(\mathrm B_{B_C}\operatorname{pr}_1)c_{o,C}.
 \end{aligned}
\]
The first remembers the same structured positive bundle.  The second
forgets the adapted filtration and retains the middle determinant
lattice induced from the same chosen $\Delta_3$ lattice.  The third
forgets the quotient generator and
its lift, retaining the original integral Tate frame.  These
forgetful identifications can be checked on the local frames just
constructed and then descend; they commute with base change and with
the specified semilinear arithmetic descent.  Applying the finite
classifying naturality and then its rational limit to
$e_3=\kappa e_J$ and \eqref{h3:eq:rat-parabolic-trace} therefore gives
$e_3=\kappa b e_2$.  Both factors are nonzero.
\end{proof}

\subsection{Witt constants and the localization map}

Define $H_W^i(Z)=H^i(R\varprojlim_\ell
R\Gamma_v(Z,W_\ell\mathcal O^\flat))$.  For $A=k$ and $A=C^\flat$,
write $\mathcal Y_{1,A}$ for $\mathcal Y_1$ and $\mathcal Y_{1,C}$,
respectively, and use $c_{o,A}$ for the corresponding ordinary
classifying map.  Give $W_\ell(A)$ its finite Witt coordinate
topology; it is discrete when $A=k$.  At each length there is a
commutative square of external derived complexes
\begin{equation}\label{h3:eq:rat-finite-witt-square}
\begin{tikzcd}[column sep=large]
 C^\bullet_{\hthreects}(G_o,\Lambda_\ell)
   \arrow[r,"\operatorname{cl}_{c_{o,A},\ell}"] \arrow[d]
 & R\Gamma_v(\mathcal Y_{1,A},\Lambda_\ell) \arrow[d]\\
 R\Gamma(G_o,\underline{W_\ell(A)})(*) \arrow[r,"\mathrm{constants}"]
 & R\Gamma_v(\mathcal Y_{1,A},W_\ell\mathcal O^\flat).
\end{tikzcd}
\end{equation}
The vertical arrows are induced by
$\Lambda_\ell=W_\ell(\F_p)\to W_\ell(A)$ and by the Witt coefficient
unit.  The bottom arrow is the actual ordinary constants map: over
$B_C$ it comes from \Cref{h3:prop:ordinary-constants}, and over $B_k$
from its base-field Frobenius descent in
\Cref{h3:thm:completed-cohomology}.  Its extension to Witt
coefficients is proved below.
On the torsor nerve both composites send a finite constant function
to the same section through
$W_\ell(\F_p)\to W_\ell(A)\to W_\ell\mathcal O^\flat$.
Thus \Cref{h3:rat:finite-classifying} proves the square term by term,
including its base change from $k$ to $C^\flat$ and its reductions in
$\ell$.  Over $B_k$, the same identity for the $P_3$ presentation
identifies the Witt image of $\operatorname{cl}_{c_{3,k},\ell}$ with
the map from $C^\bullet_{\hthreects}(P_3,\Lambda_\ell)$ to the
ordinary perfected Witt cochains in \Cref{h3:des:witt-constants}.
Here the canonical interchange of the two frame factors identifies
$G_o$ with the group $G_{\mathrm f}$ used there; the earlier geometric
maps written $BG_{\mathrm f}$ are these relative classifying maps on
their displayed bases.
Those proofs retain profinite parameters through the identification
of this quotient nerve with the computing cochains.

Taking the derived coefficient limits of the ordinary constants maps
gives the following square; their finite equivalences are justified
in the proof below:
\begin{equation}\label{h3:eq:rat-witt-square}
\begin{tikzcd}[column sep=large]
 H^i(G_o,W(k)) \arrow[r,"\sim"] \arrow[d]
     & H_W^i(\mathcal Y_1) \arrow[d]\\
 H^i(G_o,W(C^\flat)) \arrow[r,"\sim"]
     & H_W^i(\mathcal Y_{1,C}).
\end{tikzcd}
\end{equation}

\begin{lemma}\label{h3:rat:witt-nonvanishing}
The right vertical map of \eqref{h3:eq:rat-witt-square} is injective
after inversion of $p$.  The Witt coefficient image of $e_3$
on $\mathcal Y_1$ is nonzero.  Its corresponding image under
the actual ordinary deformation-coefficient map is nonzero.
\end{lemma}

\begin{proof}
For completeness, the Witt extension of the constants calculation
uses the natural additive exact sequences
\[
 0\longrightarrow R\xrightarrow{V^{\ell-1}}W_\ell(R)
    \longrightarrow W_{\ell-1}(R)\longrightarrow0.
\]
They are exact for the solid coefficient objects and v-sheaves
in question: in Witt coordinates, restriction has the continuous
set-theoretic section that appends zero.  Apply the sequence to
both sides of the characteristic-$p$ constants equivalence.
Induction using the resulting morphisms of exact triangles
proves the equivalence at every length.  Take derived inverse
limits and then the residual group invariants.  All maps in this
construction are induced by constants, so it proves the square
with its asserted arrows and their compatibility with cup products.

The group $G_o$ has no $p$-torsion.  Choose its resolution in
\Cref{h3:rat:solid-resolution} and put
$Q_\bullet=\Z_p\otimes_{\Z_p[[G_o]]}P_\bullet$.
Each $Q_j$ is finite projective, hence finite free over $\Z_p$.
For a trivial coefficient module $M$, the computing complex is
$\Hom_{\Z_p}(Q_\bullet,M)$, a fixed bounded complex of finite
matrices over $\Z_p$ tensored with $M$.
Since $W(C^\flat)$ is torsion free over the discrete valuation
ring $W(k)$, it is flat, and its fraction field extension is
faithfully flat.  Consequently
\[
 H^i(G_o,W(k))[1/p]\otimes_{W(k)[1/p]}W(C^\flat)[1/p]
       \simeq H^i(G_o,W(C^\flat))[1/p].
\]
This proves the injection in the square.  It also proves that
the image of $e_2$ is nonzero: it is the scalar extension of
the nonzero trace class of the first factor of $G_o$.

By \eqref{h3:eq:rat-finite-witt-square} and its $P_3$ counterpart,
the two Witt images of the group classes on $\mathcal Y_{1,C}$ are
the coefficient-unit images of the exact classifying images in
\eqref{h3:eq:rat-ordinary-traces}.  Their maps over $k$ base change
to these maps before the derived limit is taken.  Applying that
limit and inversion of $p$ to the geometric equality, and then
using the injection in \eqref{h3:eq:rat-witt-square}, proves the
same equality over $k$ in Witt cohomology.  Its right side is
nonzero, as just shown.

It remains to follow this nonzero Witt class through the actual
deformation-coefficient map.  Put
\[
 \mathsf C^W_\ell=C^\bullet_{\hthreects}(P_3,W_\ell B_0),\qquad
 \mathsf C^D_\ell=C^\bullet_{\hthreects}(P_3,D_\ell),\qquad
 D_\ell=(W(k)/p^\ell)[[x,y]][x^{-1}].
\]
By \Cref{h3:lem:ghost-comparison}, the two maps from $W_\ell B_0$
to $D_\ell$ and to $W_\ell B_{\hthreepk}$ induce cohomology
isomorphisms.  The graded maps for $\gamma_\ell$ are powers of
$\phi_p$, while those for $\beta_\ell$ are the natural maps
$B_0\to B_{\hthreepk}$.  Both are cohomological equivalences by
\Cref{h3:thm:ordinary-comparison}.  This statement is used at
each finite length; the two inverse systems have different
transition maps.

Choose an integral constant class
$\beta\in H^3_{\mathrm{int}}(P_3)$ with
$\beta[1/p]=p^N e_3$ for some $N\geq0$, where $e_3$ denotes the
source class in $H^3_{\mathrm b}(P_3)$.  The normalized rational
trace is fixed by $\Gamma=\Gal(k/\F_p)$, as proved in
\Cref{h3:rat:trace-groups}.  Thus each difference
$\gamma\beta-\beta$ is killed by a power of $p$.  Since $\Gamma$
is finite, one common power kills all these differences.  Replacing
$\beta$ by that multiple and increasing $N$ makes the integral
cohomology class $\Gamma$-invariant while retaining the same nonzero
rational trace line.  Let $\bar\beta_\ell$ be its canonical projection
to $H^3(C^\bullet_{\hthreects}(P_3,\Lambda_\ell))$.  These finite
cohomology classes are reduction-compatible and $\Gamma$-invariant.
Let
$a\in H^3(R\varprojlim_{R_\ell}\mathsf C^W_\ell)$ be its coefficient image,
where $R_\ell:W_\ell B_0\to W_{\ell-1}B_0$ is ordinary Witt restriction.
The maps to perfected and geometric Witt coefficients commute with
these restrictions.  The nonvanishing already proved therefore implies
$a[1/p]\ne0$.  The groups $H^2(\mathsf C^W_\ell)$ are finite, by
the finite-length additive Witt filtration and
\Cref{h3:prop:bounded-pole-finiteness,h3:thm:ordinary-comparison}.
Their tower is Mittag--Leffler, so the Milnor sequence gives
\[
 H^3(R\varprojlim_{R_\ell}\mathsf C^W_\ell)
      =\varprojlim_{R_\ell} H^3(\mathsf C^W_\ell).
\]
Consequently the component classes $a_\ell$ have unbounded
$p$-power orders.  By naturality of the projections, $a_\ell$ is the
coefficient image of $\bar\beta_\ell$ in $W_\ell B_0$ cohomology.

Write $d_\ell$ for their ghost images in $H^3(\mathsf C^D_\ell)$.
These have exactly the same orders.  The ghost maps fix
$\Z/p^\ell$, so $d_\ell$ is the direct coefficient image
of $\bar\beta_\ell$ and is compatible with coefficient reduction.
More explicitly the transition identity is
\[
 \operatorname{red}_{\ell,\ell-1}\gamma_\ell
       =\gamma_{\ell-1}R_\ell W_\ell(\phi_p),
\]
where $\phi_p(b)=b^p$ is the characteristic-$p$ Frobenius defined
in \Cref{h3:lem:ghost-comparison}.  The class $a_\ell$ is fixed by
$W_\ell(\phi_p)$ because it comes from $W_\ell(\F_p)$.
The naturalities of the finite units also make each $d_\ell$
$\Gamma$-invariant.  Exact semilinear descent in
\Cref{h3:lem:framework-semilinear,h3:prop:exact-descent} identifies, by restriction,
\[
 H^3_{\hthreects}(G_3(k),D_\ell)
       \simeq H^3_{\hthreects}(P_3,D_\ell)^\Gamma.
\]
Here descent is applied after extension to the semilinear $D_\ell$
coefficients.
The descended class restricts to $d_\ell$, so it has the same order;
these identifications also preserve ordinary reduction.  Thus the
specified descended constant classes have unbounded $p$-power order.
The natural
spectral comparison in \Cref{h3:des:minus-three-edge,h3:prop:integral-delta}, with its
unique bidegree $(s,t)=(3,0)$ in stem $-3$, identifies this family
with the finite projections of a class in $\pi_{-3}L_{K(1)}T$.
Its unbounded finite orders make that class nonzero after
rationalization, and the identification preserves the actual maps
from constants.
\end{proof}

\begin{lemma}[Origin in the rational sphere]\label{h3:lem:constants-origin}
The reduced-trace class in $H^3_{\mathrm b}(P_3)$, after the
constant coefficient map and coefficient-field descent, is the
degree-three class of an element
$\alpha_3\in\pi_{-3}L_0T$.  This element is nonzero and its
image in $\pi_{-3}L_0L_{K(1)}T$ is nonzero.
\end{lemma}

\begin{proof}
Reduction followed by the ordinary coefficient localization
$E_{3,0}(k)\to D_\ell$ sends the integral constant class chosen in
\Cref{h3:rat:witt-nonvanishing} to its specified finite classes
$d_\ell$.  These have unbounded $p$-power order and are compatible
with ordinary reduction.  Hence the coefficient image of the
rational trace in $H^3(P_3,E_{3,0}(k))[1/p]$ is nonzero.
The trace construction and its coefficient map commute with the
semilinear finite Galois action; their class is invariant and
descends by \Cref{h3:lem:framework-semilinear}.

We explain why this cohomology class is an actual rational
homotopy class.  A central element $z=1+p$ acts trivially on
$E_{3,0}(k)$ and by the nontrivial scalar $z^{-q}$ on the
nonzero periodicity twist $E_{3,2q}(k)$, for $q\ne0$.
This is the action on $\omega^{\otimes q}$ in the
contravariant deformation convention of
\Cref{h3:lem:invariant-differential}.  A central group element acts as the
identity on group cohomology when its coefficient action is
combined with its trivial inner conjugation.  Equivalently,
the natural transformation furnished by that element gives
a homotopy to the identity on the bar resolution.  Thus
$z^{-q}-1$ annihilates
$H^s(P_3,E_{3,2q}(k))$.  It is a nonzero element of $\Z_p$,
so every such row vanishes after inversion of $p$.

The descent of \Cref{h3:prop:exact-descent} is exact and strongly
convergent.  The finite stabilizer resolutions and exact semilinear
descent of \Cref{h3:rat:solid-resolution,h3:lem:framework-semilinear}
bound its cohomological amplitude.
Rationalization therefore preserves its finite filtrations.
Only internal degree zero remains.  After finite coefficient-field
descent, the actual coefficient maps fit into the commutative square
\[
\begin{tikzcd}[column sep=2em]
 H^3_{\hthreects}(G_3(k),E_{3,0}(k))[1/p]
   \arrow[r] \arrow[d,"\simeq"]
 & \bigl(\varprojlim_\ell
       H^3_{\hthreects}(G_3(k),D_\ell)\bigr)[1/p]
   \arrow[d,"\simeq"] \\
 \pi_{-3}L_0T \arrow[r]
 & \pi_{-3}L_0\operatorname*{holim}_\ell L_1(T/p^\ell).
\end{tikzcd}
\]
The upper arrow is induced by the coefficient reductions and
localizations just described.  The left isomorphism is the
single-row rational descent calculation.  The right isomorphism
uses the natural, unique $(s,t)=(3,0)$ edge of
\Cref{h3:des:finite-length-model} and the vanishing of the Milnor
$\lim^1$ term in \Cref{h3:prop:integral-delta}.  The finite-length
model identifies those coefficient maps with the actual Moore
quotient and localization maps on every descent term.
For commutativity, first choose an integral representative of a
$p$-power multiple of a rational source cohomology class.
Finite source filtration and rational collapse allow a further
$p$-power multiple to survive to homotopy.  Naturality in
\Cref{h3:prop:exact-descent} identifies its finite target edges
with that same multiple of the coefficient images, for every
$\ell$.  This multiplier is chosen in the source, independently
of $\ell$.  Taking the coefficient limit and then inverting $p$
gives the displayed square.

By \Cref{h3:lem:height-one-completion}, the bottom map is the
rationalization of $T\to L_{K(1)}T$.  The nonzero rational trace
image in the upper left therefore determines
$\alpha_3\in\pi_{-3}L_0T$, and its image in the upper right is the
nonzero rational class of the compatible finite constants.
The square proves the asserted nonvanishing of its actual
localization without an additional arithmetic coefficient splitting.
\end{proof}

\begin{theorem}[Rational compatibility]\label{h3:thm:rational-compatibility}
The generators in \eqref{h3:eq:framework-rational-boundary} may be
chosen so that $\alpha_1=\zeta$ and the actual localization map
$T\to L_{K(1)}T$ satisfies
\begin{equation}\label{h3:eq:rat-actual-images}
 \alpha_3\longmapsto c\delta,\qquad
 \zeta\alpha_3\longmapsto c\zeta\delta
       \qquad(c\in\Q_p^\times)
\end{equation}
after rationalization.  In particular the rational homotopy of
$T$ has basis
\[
 1,\ \zeta,\ \alpha_3,\ \zeta\alpha_3,\
 \alpha_5,\ \zeta\alpha_5,\ \alpha_3\alpha_5,\
 \zeta\alpha_3\alpha_5
\]
in degrees $0,-1,-3,-4,-5,-6,-8,-9$, respectively.
\end{theorem}

\begin{proof}
The rational exterior calculation
\eqref{h3:eq:framework-rational-boundary} makes degree $-3$
one-dimensional.  Choose its generator to be the nonzero
constant-origin class of \Cref{h3:lem:constants-origin}.  The
determinant class is the degree-one group primitive by its
construction in \Cref{h3:sec:descent}, and its actual image is
nonzero rationally by \Cref{h3:thm:height-one-maps}.  Thus it
may be used as $\alpha_1$.  Choose any remaining exterior generator in
degree $-5$.

By \Cref{h3:prop:integral-delta}, the rational target in degree
$-3$ is precisely $\Q_p\delta$.  Nonvanishing of the actual
image gives $\alpha_3\mapsto c\delta$ with $c\ne0$.
All coefficient comparisons and localizations used above
respect products, and the same actual determinant class
occurs on source and target.  Multiplication by it gives
the second formula with the same scalar.  The exterior
calculation then supplies the displayed basis.
\end{proof}
